\subsection{How to tell whether a milestone has been met}
\label{sec:25.1.2}
A milestone claimed by the organization that built the system is a different kind of evidence from one checked by somebody else, and the gap is empirical rather than about trust. A system's own claim inherits the self-report limits; a passing score on a fixed suite inherits specification gaming and goal misgeneralization. What either needs instead is a passing score under conditions built to break it.

Who runs it matters as much as the test. METR evaluates both in partnership with developers and separately after release \autocite{metr2026about}. The UK AI Security Institute does government-run pre-deployment evaluation and is independent — and by its own account is not a regulator: it can't compel submission, block a release, or intervene once a system is causing harm \autocite{aisi2025governance}. Independent evaluation and enforcement are not the same capability.

An evaluation like this is worth running only if it is willing to report that a milestone wasn't met — and, where the milestone measured the wrong thing, revising the milestone. A negative result may mean the work fell short or that the milestone was wrong, and evaluators have to distinguish those before deciding what to revise.
